\section{Closing}
\label{sec:conclusions_closing}

This thesis has shown that computational, data-driven approaches can support both the design and the operation of docked \acp{bss}. In Part~\ref{pt:data-driven-station-planning}, a flexible station-planning framework was developed and applied to Barcelona, enabling multi-criteria, network-aware, and topography-sensitive placement and supporting scenario exploration and incremental expansion. In Part~\ref{pt:predictive_operations}, the spatiotemporal mobility patterns of \textit{Bicing} were documented, a joint framework for \ac{e-b} locations and battery levels was then presented, and interpretable \ac{pm} models for key components of both \acp{m-b} and \acp{e-b} were built and validated. Applied to Barcelona and \textit{Bicing}, these methods illustrate how operational and open urban data can be turned into actionable insights for planners and operators. 

These contributions are methodologically transferable, as they rely on the kind of data that many \acp{bss} and cities already collect or can obtain from open sources. They are therefore aligned with the growing availability of high-resolution urban and \ac{bss} data, and can support more efficient, equitable, and sustainable shared mobility as systems mature and data practices expand.
